% Section 6, part 9.
Let $c$ be a branch point, $t:=p-c$, and $\chi$ (a local letter) a local coordinate at $c$ with $x=x(c)+\lambda\chi^2$, $\lambda\neq 0$, so that $\iota_c\colon\chi\mapsto-\chi$. Write $y=\sum_ky_k\chi^k$ and, near $(c,c)$, $\omega_{0,2}=\big[(\chi_1-\chi_2)^{-2}+\mathrm{reg}(\chi_1,\chi_2)\big]\dd\chi_1\dd\chi_2$ with $\mathrm{reg}$ holomorphic and symmetric; then $\mathrm{reg}_{00}:=\mathrm{reg}(0,0)=\frac16\{t;\chi\}(0)$, where $\{t;\chi\}=t'''/t'-\frac32(t''/t')^2$ is the Schwarzian derivative of $t$ as a function of $\chi$ (both sides are $t_3/t_1-(t_2/t_1)^2$ for $t=t_1\chi+t_2\chi^2+t_3\chi^3+\cdots$, by expanding $t'(\chi_1)t'(\chi_2)/(t(\chi_1)-t(\chi_2))^2$). For a $1$-form $\eta=\sum_m\eta_m\chi^m\dd\chi$ holomorphic at $c$ put
\begin{equation}
\label{eq:mod-loc}
\mathrm{Loc}_c[\eta]:=-\frac{1}{4\lambda y_1}\Big[\mathrm{reg}_{00}\,\eta_0+\frac{\eta_2}{12}-\frac{y_3}{4y_1}\,\eta_0\Big].
\end{equation}

\begin{lem}[A local formula for $\omega_{1,1}$]\genusstatus{AI-proposed, unchecked}
\label{lem:mod-w11}
\begin{enumerate}
\item $\omega_{1,1}(p_0)=\sum_c\mathrm{Loc}_c[\omega_{0,2}(p_0,\cdot)]$ for $p_0$ off the branch points, $\omega_{0,2}(p_0,\cdot)$ being a form in the second slot.
\item $\omega_{1,1}=\dd G$ with $G$ rational, with poles only at the branch points and $G(\infty)=0$; so $\omega_{1,1}$ has zero residues, is holomorphic at $p=0$ and is $O(\dd p/p^2)$ at $\infty$.
\item $\langle-\log p\rangle^{\mathrm{TR}}_1=\int_{\mathfrak r}\omega_{1,1}=\sum_c\mathrm{Loc}_c[-\dd q/q]$.
\end{enumerate}
\end{lem}

\begin{proof}
(i) By \eqref{eq:mod-kernel},
\[
\omega_{1,1}(p_0)=\sum_c\Res_{q\to c}K_c(p_0,q)\,\omega_{0,2}(q,\iota_cq).
\]
Expand $\omega_{0,2}(p_0,q')=\dd p_0\sum_mB_m(p_0)\chi^m\dd\chi$ for $q'$ near $c$ with coordinate $\chi$. Then $\int_{\iota_cq}^q\omega_{0,2}(p_0,\cdot)=2\,\dd p_0\sum_{m\text{ even}}B_m\chi^{m+1}/(m+1)$, $y(q)-y(\iota_cq)=2\sum_{k\text{ odd}}y_k\chi^k$ and $\dd x=2\lambda\chi\,\dd\chi$, so
\[
K_c(p_0,q)=\frac{\dd p_0}{\dd\chi}\,\frac{1}{4\lambda y_1}\Big[B_0\chi^{-1}+\Big(\frac{B_2}{3}-B_0\frac{y_3}{y_1}\Big)\chi+O(\chi^3)\Big],
\]
and $\omega_{0,2}(q,\iota_cq)=-\big[(4\chi^2)^{-1}+\mathrm{reg}_{00}+O(\chi^2)\big]\dd\chi^2$ ($\dd(-\chi)=-\dd\chi$, and $\mathrm{reg}(\chi,-\chi)=\mathrm{reg}_{00}+O(\chi^2)$ by symmetry). The coefficient of $\chi^{-1}\dd\chi$ in the product is $-\frac{1}{4\lambda y_1}\big[B_0\,\mathrm{reg}_{00}+\frac14\big(\frac{B_2}{3}-B_0\frac{y_3}{y_1}\big)\big]\dd p_0=\mathrm{Loc}_c[\omega_{0,2}(p_0,\cdot)]$.
(ii) $\omega_{0,2}(p_0,q)=\dd_{p_0}\big[\dd q/(q-p_0)\big]$, and near $c$, $\dd q/(q-p_0)=\sum_mg_m(p_0)\chi^m\dd\chi$ with $g_m$ rational, with poles only at $c$ and $g_m(\infty)=0$. So $B_m\dd p_0=\dd g_m$, and by (i) $\omega_{1,1}=\dd G$ with $G$ a combination of the $g_0,g_2$ at the five branch points. (iii) The residues in (i) are integrals over small circles avoiding $\mathfrak r$; the ray integral converges by (ii), and as $\omega_{0,2}(p_0,q)=O(p_0^{-2})\dd p_0$ uniformly for $q$ on the circles, it may be taken inside: $\int_{p_0\in\mathfrak r}\omega_{0,2}(p_0,q)=-\dd q/q$, and $\mathrm{Loc}_c$ is linear in finitely many Taylor coefficients. Finally, by (i) and \eqref{eq:mod-kernel}, $\mathrm{Loc}_c[\eta]=\Res_{q\to c}\frac12\frac{\int_{\iota_cq}^q\eta}{(y(q)-y(\iota_cq))\dd x(q)}\,\omega_{0,2}(q,\iota_cq)$ for $\eta$ holomorphic at $c$; this form does not depend on the coordinate $\chi$.
\end{proof}

\begin{thm}[Theorem V]\genusstatus{AI-proposed, unchecked}
\label{thm:mod-V}
Let $h\geq 0$, $a,k\geq 0$, $n=a+k\geq 1$, $2h-2+n>0$, and $\phi_1,\dots,\phi_a\in\bbQ[v][p]$.
\begin{enumerate}
\item If $a\geq 1$, then $\langle\phi_1\cdots\phi_a,(-\log p)^k\rangle^{\mathrm{TR}}_h\in\bbQ[v]$.
\item If $a=0$, then $\langle(-\log p)^k\rangle^{\mathrm{TR}}_h=(\FTR_h)^{(k)}(v)$, with $\FTR_h$ of \cref{thm:mod-gap} for $h\geq 2$, $\FTR_1:=-\frac1{12}\log v$ (here $k\geq 1$) and $\FTR_0:=\frac12v^2\log v$ (here $k\geq 3$).
\item The recursion \eqref{eq:mod-V1} holds.
\end{enumerate}
\end{thm}

\begin{proof}
(iii) is \cref{prop:mod-V1}. (i) Fix $h$ and $a\geq 1$ and induct on $k$. The first admissible case is $k=0$ if $2h-2+a>0$ (\cref{thm:mod-Lreg}), and otherwise $\langle\phi_1,(-\log p)^2\rangle^{\mathrm{TR}}_0$ or $\langle\phi_1\phi_2,-\log p\rangle^{\mathrm{TR}}_0$ (\cref{lem:mod-genus0}~(iv)). In the step, the right side of \eqref{eq:mod-V1} consists of a $v$-derivative of, and pairings with, $a$ L-legs in $\bbQ[v][p]$ ($p^5-v$ is monic) and $k$ log legs; all lie in $\bbQ[v]$ by induction.

(ii) For $a=0$, \eqref{eq:mod-V1} reads $\langle(-\log p)^{k+1}\rangle^{\mathrm{TR}}_h=\partial_v\langle(-\log p)^k\rangle^{\mathrm{TR}}_h$, so it suffices to treat the first admissible $k$. \emph{$h\geq 2$, $k=1$.} As in the proof of \cref{prop:mod-period}~(iii), $(2-2h)\FTR_h=v\langle-\log p\rangle^{\mathrm{TR}}_h$, and $\FTR_h$ is a multiple of $v^{2-2h}$ (\cref{thm:mod-gap}), so $\langle-\log p\rangle^{\mathrm{TR}}_h=(\FTR_h)'(v)$. \emph{$h=0$, $k=3$.} \Cref{lem:mod-genus0}~(iv). \emph{$h=1$, $k=1$.} We show $\langle-\log p\rangle^{\mathrm{TR}}_1=-1/(12v)$.

\emph{Reduction to $v=1$.} By \cref{lem:mod-scale}~(ii) with $k=\varsigma$, $m_\varsigma^*\omega^{(v)}_{1,1}=v^{-1}\omega^{(1)}_{1,1}$, and rays go to rays, so $\langle-\log p\rangle^{\mathrm{TR}}_1(v)=v^{-1}\langle-\log p\rangle^{\mathrm{TR}}_1(1)$. Let $v=1$, $\varsigma=1$, $\varsigma_b=\zeta^{-b}$.

\emph{Five equal contributions.} $\rho_b(p):=\zeta^{-b}p$ sends $1$ to $\varsigma_b$, with $x\circ\rho_b=\zeta^bx$, $y\circ\rho_b=\zeta^{-b}y$, $\rho_b^*\omega_{0,2}=\omega_{0,2}$, $\rho_b^*(\dd q/q)=\dd q/q$; hence $\iota_b\circ\rho_b=\rho_b\circ\iota_0$ and $(y(q)-y(\iota q))\,\dd x(q)$ is unchanged. By the coordinate-free form of $\mathrm{Loc}$, $\mathrm{Loc}_{\varsigma_b}[-\dd q/q]=\mathrm{Loc}_1[-\dd q/q]$.

\emph{The contribution at $c=1$.} With $t=p-1$, $x(1+t)=\frac54+\frac52t^2+\frac54t^4-t^5+O(t^6)$, so $\lambda=5/2$, $\chi=t+\frac{t^3}4-\frac{t^4}5+O(t^5)$ and $t=\chi-\frac{\chi^3}4+\frac{\chi^4}5+O(\chi^5)$. Hence $y_1=1$, $y_2=0$, $y_3=-\frac14$; $\{t;\chi\}(0)=-\frac32$ ($t'''(0)=-\frac32$, $t''(0)=0$), so $\mathrm{reg}_{00}=-\frac14$; and $-\dd q/q=-\big(1-\chi+\frac14\chi^2+O(\chi^3)\big)\dd\chi$, so $\eta_0=-1$, $\eta_2=-\frac14$. By \eqref{eq:mod-loc}, $\mathrm{Loc}_1[-\dd q/q]=-\frac1{10}\big[\frac14-\frac1{48}-\frac1{16}\big]=-\frac1{60}$. By \cref{lem:mod-w11}~(iii) the total at $v=1$ is $5\cdot(-\frac1{60})=-\frac1{12}$. So $\langle-\log p\rangle^{\mathrm{TR}}_1=-1/(12v)=(\FTR_1)'(v)$.
\end{proof}

% Section 6, part 10.
\subsection{The Model Lemma on \texorpdfstring{$\sK$}{K}}
\label{sec:mod-MLK}

Recall $\ML{h}$ from \cref{def:red-ML}, with $\Fm_h$ the free energy of \cref{cor:mod-trF} for $h\geq 2$, $\Fm_1=-\frac1{12}\log v$ and $\Fm_0=-\frac52v^2\log v$ (only their derivatives of order $\geq 1$, respectively $\geq 3$, enter). Then $\Fm_h=(-5)^{1-h}\FTR_h$ for every $h\geq 0$, with $\FTR_0,\FTR_1$ as in \cref{thm:mod-V}~(ii).

\begin{thm}[The Model Lemma on $\sK$]\genusstatus{AI-proposed, unchecked}
\label{thm:mod-MLK}
$\ML{h}$ holds for every $h\geq 0$.
\end{thm}

\begin{proof}
Part (ii) of $\ML{h}$ is \cref{thm:mod-gap}. For part (i), let $X_i\in\sK\cap V_{w_i}$, $i=1,\dots,n$.

\emph{Step 1 (the legs of $\sK$).} By \cref{cor:mod-Ksplit}, $X_i=Y_i+\Pi_0(X_i)H$ with $Y_i\in\sKo\cap V_{w_i}$ and $\Pi_0(X_i)\in\bbQ v^{w_i/5}$, a monomial with a non-negative exponent. By \cref{prop:mod-LD} ($q_{-\log p}=H$, and the TR leg of $Y_i$ is $-\Lambda(Y_i)$), $X_i=q_{\ell_i}$ for $\ell_i:=-\Lambda(Y_i)+\Pi_0(X_i)\,(-\log p)$. By \cref{thm:mod-lattice}~(b) and \cref{lem:mod-classL}~(i), $-\Lambda(Y_i)=\phi_i+f_i(x)$ with $\phi_i\in\bbQ[v][p]$ and $f_i\in\bbQ[v][x]$.

\emph{Step 2 (expansion).} \Cref{thm:mod-tr,thm:mod-Lreg,thm:mod-V} are stated for a complex number $v\neq0$, and $\ML{h}$ for the variable $v$. This makes no difference: the correlator $\langle X_1\cdots X_n\rangle_h$ is a finite $\bbQ[v]$-combination of monomials $Cv^e$, $C\in\bbQ$, $e\in\bbZ$ (\cref{lem:set-modhom}), every term on the right below is the value at $v$ of an element of $\bbQ[v,v^{-1}]$, and an identity between elements of $\bbQ[v,v^{-1}]$ that holds for all complex $v\neq0$ holds identically. (The derivatives $(\Fm_h)^{(n)}$ that occur, of order $\geq1$ for $h=1$ and $\geq3$ for $h=0$, are in $\bbQ[v,v^{-1}]$.) By \cref{thm:mod-tr}, multilinearity and \cref{lem:mod-xfun}, which removes the components $f_i(x)$ in every slot,
\[
\langle X_1\cdots X_n\rangle_h=(-5)^{1-h}\sum_{S\subseteq\{1,\dots,n\}}\prod_{i\notin S}\Pi_0(X_i)\,\big\langle(\phi_i)_{i\in S},(-\log p)^{n-|S|}\big\rangle^{\mathrm{TR}}_h .
\]
Every term with $S\neq\emptyset$ lies in $\bbQ[v]$, by \cref{thm:mod-Lreg} if $S=\{1,\dots,n\}$ and by \cref{thm:mod-V}~(i) otherwise.

\emph{Step 3 (the pure logarithmic term).} By \cref{thm:mod-V}~(ii) and $\Fm_h=(-5)^{1-h}\FTR_h$, the term $S=\emptyset$ is $(-5)^{1-h}\langle(-\log p)^n\rangle^{\mathrm{TR}}_h\prod_i\Pi_0(X_i)=(\Fm_h)^{(n)}\prod_i\Pi_0(X_i)$.
\end{proof}

\begin{cor}[The one-point law on $\sK$]\genusstatus{AI-proposed, unchecked}
\label{cor:mod-onepoint}
Let $h\geq 1$ and $X\in\sK\cap V_w$. If the $v$-exponent $e=1-2h+w/5$ of $\langle X\rangle_h$ (\cref{lem:set-modhom}) is negative, then $\langle X\rangle_h=\Pi_0(X)\,(\Fm_h)'(v)$.
\end{cor}

\begin{proof}
By \cref{thm:mod-MLK} with $n=1$, $\langle X\rangle_h-(\Fm_h)'\Pi_0(X)\in\bbQ[v]$. By \cref{lem:set-modhom} both terms are multiples of $v^e$, $e=1-2h+w/5<0$ (as $\Pi_0(X)\in\bbQ v^{w/5}$), so the difference is $0$.
\end{proof}

\begin{rmk}[Off $\sK$]
\label{rmk:mod-offK}
The restriction to $\sK$ cannot be dropped: for insertions with a component along some $\psi^kH$, $k\geq 1$, the singular part of a one-point correlator is in general not $\Pi_0(X)(\Fm_h)'$; this is known only by computation (\cref{sec:evi-model}) and are not used.
\end{rmk}

